
\documentclass[12pt]{article}

% Layout.
\usepackage[top=1in, bottom=0.75in, left=1in, right=1in, headheight=1in, headsep=6pt]{geometry}

% Fonts.
\usepackage{mathptmx}
\usepackage[scaled=0.86]{helvet}
\renewcommand{\emph}[1]{\textsf{\textbf{#1}}}

% TiKZ.
\usepackage{tikz, pgfplots}
\usetikzlibrary{calc}
\pgfplotsset{compat = newest}
 
\pgfplotsset{my style/.append style={axis x line=middle, axis y line=
middle, xlabel={\(x\)}, ylabel={\(y\)}, axis equal }}

% Misc packages.
\usepackage{amsmath,amssymb,latexsym}
\usepackage{graphicx}
\usepackage{array}
\usepackage{xcolor}
\usepackage{multicol}

% Commands to set various header/footer components.
\makeatletter
\def\doctitle#1{\gdef\@doctitle{#1}}
\doctitle{Use {\tt\textbackslash doctitle\{MY LABEL\}}.}
\def\docdate#1{\gdef\@docdate{#1}}
\docdate{Use {\tt\textbackslash docdate\{MY DATE\}}.}
\def\doccourse#1{\gdef\@doccourse{#1}}
\let\@doccourse\@empty
\def\docscoring#1{\gdef\@docscoring{#1}}
\let\@docscoring\@empty
\def\docversion#1{\gdef\@docversion{#1}}
\let\@docversion\@empty
\makeatother

% Headers and footers layout.
\makeatletter
\usepackage{fancyhdr}
\pagestyle{fancy}
\fancyhf{} % Clears all headers/footers.
\lhead{\baselineskip 30pt
%\emph{\@doctitle\hfill\@docdate}
\emph{\@docdate\hfill\@doctitle}
\ifnum \value{page} > 1\relax\else\\
\emph{Name: \rule{3.5in}{1pt}\hfill \@docscoring}\fi}
\rfoot{\emph{\@docversion}}
\lfoot{\emph{\@doccourse}}
\cfoot{\emph{\thepage}}
\renewcommand{\headrulewidth}{0pt}%
\makeatother

% Paragraph spacing
\parindent 0pt
\parskip 6pt plus 1pt

% A problem is a section-like command. Use \problem{5} to
% start a problem worth 5 points.
\newcounter{probcount}
\newcounter{subprobcount}
\setcounter{probcount}{0}
\newcommand{\problem}[1]{%
\par
\addvspace{4pt}%
\setcounter{subprobcount}{0}%
\stepcounter{probcount}%
\makebox[0pt][r]{\emph{\arabic{probcount}.}\hskip1ex}\emph{[#1 points]}\hskip1ex}
\newcommand{\thesubproblem}{\emph{\alph{subprobcount}.}}

% Subproblems are an enumerate-like environment with a consistent
% numbering scheme. 
% Use \begin{subproblems}\item...\item...\end{subproblems}
\newenvironment{subproblems}{%
\begin{enumerate}%
\setcounter{enumi}{\value{subprobcount}}%
\renewcommand{\theenumi}{\emph{\alph{enumi}}}}%
{\setcounter{subprobcount}{\value{enumi}}\end{enumerate}}

% Blanks for answers in normal and math mode.
\newcommand{\blank}[1]{\rule{#1}{0.75pt}}
\newcommand{\mblank}[1]{\underline{\hspace{#1}}}
\def\emptybox(#1,#2){\framebox{\parbox[c][#2]{#1}{\rule{0pt}{0pt}}}}

% Misc.
\renewcommand{\d}{\displaystyle}
\newcommand{\ds}{\displaystyle}
\def\bc{\begin{center}}
\def\ec{\end{center}}
\def\be{\begin{enumerate}}
\def\ee{\end{enumerate}}


% Document info
\doctitle{Math F251X: Quiz }
\docdate{April 2, 2026}
\doccourse{UAF Calculus I}
\docversion{Spring 2026}
\docscoring{\blank{0.8in} / 25}

\begin{document}

\emph{Please circle your instructor's name:}
\hfill James Gossell \hfill Gordon Williams \\

There are 25 points possible on this quiz.
Any outside materials are not allowed.
\emph{For full credit, show all work clearly.}


\problem{10} Find the radius, \(r\) , and the height, \(h\), of the closed cylinder with volume of \(16\pi\) that has the least amount of surface area. The formulas for the volume, \(V\), and the surface area, \(S\), are given below.
\[V=\pi r^{2}h, \hspace{1cm} S=2\pi rh+2\pi r^{2}\]
\includegraphics[width=1.5in]{Quiz9Cylinder}
\vfill
\problem{6} Solve the initial value problems given below.
%499,501,505,507

\begin{subproblems}
\item Suppose $f'(x)=3x^{2}-1$, and $f(1)=5$. What is $f(x)$?
\vspace{1in}
\item Suppose $f'(x)=3\sin(x)+e^{-x}$ and $f(0)=0$. What is $f(x)$?
\vspace{1in}
\end{subproblems}

\newpage
\problem{9}  Evaluate the following limits. \emph{Show your work}, including appropriate use of limit notation. If you use L'H\^opital's rule, you must indicate where you are using it by writing \(\stackrel{H}{=}\) or \(\stackrel{L'H}{=}\) or something similar. Use \(\infty\) or \(-\infty\) where appropriate, and if the limit does not exist, write {\sf DNE} and provide a justification.

\begin{subproblems}%#358, 360, 367, 368, 369, 371, 372, 373, 377, 380, 382, 385, 387,395
\item \(\displaystyle \lim_{x\to 0} \dfrac{(1+x)^{1/3}-1-\dfrac{1}{3}x}{x^{2}}\) %Like 374 double application of LH
\vfill
\item \(\displaystyle \lim_{x\to \frac{\pi}{4}}(1-\cos(x))\csc(x) \) %Like 390. DOES NOT INVOLVE L'H.
\vfill
\item \(\displaystyle \lim_{x\to \infty}\left(1+\dfrac{3}{x}\right)^{x} \)% Involves logarithmic method.
\vfill
\vfill
\end{subproblems}




\end{document}